\section{Effect of Batch Size on Latency and Cost (Stage 1)}
\label{sec:effect_of_batch_size}

Stage 1 sweeps evaluated the performance of a Fixed Batching Policy under varying constraints of maximum batch size. Analysis of the raw data reveals the fundamental latency-cost trade-off inherent in rollup design: larger batch sizes reduce per-transaction L1 gas costs by amortizing fixed submission overhead, but exponentially increase mempool queuing delays.

For instance, configuring the sequencer with a batch size of 25 resulted in a low average queue wait time of 1,392 ms, but incurred a relatively high gas cost of approximately 20,549 gas/tx. Conversely, increasing the batch size limit to 500 drastically reduced the amortized gas cost to roughly 19,459 gas/tx, but at the severe penalty of inflating the queue wait time to 74,209 ms. 

It is critical to note that the raw Stage 1 metrics exhibited boundary-effect distortions. Specifically, empty tail batches submitted at the end of finite test runs artificially inflated the simple average gas/tx for large configurations. Thus, evaluating the steady-state performance must account for this tail distortion.

\section{Cost Amortization and Distortions}
\label{sec:cost_results}
Rollup L1 publication costs exhibit dual-layer characteristics: a fixed cost associated with submitting the batch and verifying the validity proof, and a variable cost that scales with transaction data size. Increasing batch sizes allows the sequencer to amortize the fixed costs over more transactions, reducing the per-transaction L1 gas fee.
In the Stage 1 sweeps (utilizing a Mock Verifier contract with a low ~29,000 gas overhead), the true weighted gas per transaction scaled as follows:
\begin{itemize}
    \item \textbf{Batch Size 25 (s1\_bs\_0025):} True Weighted Gas/Tx was \textbf{20,549 gas}.
    \item \textbf{Batch Size 100 (s1\_bs\_0100):} True Weighted Gas/Tx was \textbf{19,681 gas}.
    \item \textbf{Batch Size 500 (s1\_bs\_0500):} True Weighted Gas/Tx was \textbf{19,459 gas}.
\end{itemize}
\subsection{Empty Batch Cost Distortion}
A significant statistical anomaly was identified in the unweighted aggregation framework. For the large batch run \texttt{s1\_bs\_1000}, the reported gas cost was \textbf{83,626 gas/tx}, whereas smaller batches reported ~19,600 gas. 
Telemetry logs showed that the run produced 6 batches: Batches 1 and 2 were full (605 and 661 txs), Batch 3 had 1 tx, and Batches 4, 5, and 6 had 0 txs (empty). Because the metrics collector uses a simple mean of the ratios and applies a floor of $\max(\text{tx\_count}, 1)$ to prevent division by zero, the empty batches were calculated as:
\begin{equation}
\text{Gas/Tx}_{\text{empty}} = \frac{112,332\text{ gas}}{1\text{ tx}} = 112,332\text{ gas/tx}
\end{equation}
Averaging these ratios yields the distorted 83.6k gas/tx. The true weighted average (total gas divided by total transactions) was actually \textbf{19,771 gas/tx}, matching the amortization curve. This underscores the necessity of using true weighted averages when reporting rollup transaction economics.

\section{Queueing Latency (Soft Finality)}
\label{sec:latency_results}
The Stage 1 experiments established a clear relationship between batch size configurations and queueing latency (the time a transaction waits in the sequencer pool before the batch is sealed).

As the timeout interval was increased (from 500 ms to 10,000 ms), the average queue wait time scaled linearly until it hit the ceiling defined by the \texttt{MAX\_BATCH\_SIZE}. For example:
\begin{itemize}
    \item \textbf{Timeout 500 ms}: Actual queue wait was 256.4 ms.
    \item \textbf{Timeout 2000 ms}: Actual queue wait was 1017.2 ms.
    \item \textbf{Timeout 10000 ms}: The batch filled to the size cap (100 txs) in roughly 5.4 seconds due to the arrival rate, resulting in an actual queue wait of 2937.1 ms instead of the full 10 seconds.
\end{itemize}

A proposed visualization for these findings is a dual-axis plot showing Timeout on the X-axis, Average Queue Wait Time on the Left Y-Axis, and Average Batch Size on the Right Y-Axis. This plot clearly illustrates how the timeout acts as the primary batch sealing trigger until the size cap is reached.
